% Passes 5611--5618: physics-first magnetic frontier.
\subsection{Heisenberg magnetic bulk, intrinsic vector lift, and finite gauge dynamics}

The projectivity-frame programme now admits a sharper separation between finite kinematics,
phase data, and candidate dynamics.  On the affine bulk $V=\mathbb F_q^2$, define
\[
 A_q((x,y),(u,v))={\bf 1}_{x\ne u,\,y\ne v}\,\chi(xv-uy),
\]
where $\chi$ is a nontrivial additive character.  Writing $K$ for the full symplectic
Fourier kernel and $R,C$ for the same-coordinate kernels gives
\[
 A=K-R-C+I,
\]
with
\[
 K^2=q^2I,\qquad R^2=qR,\qquad C^2=qC,
\]
\[
 KR=RK=qR,\qquad KC=CK=qC,\qquad RCR=qR,\qquad CRC=qC.
\]
Hence the exact affine magnetic spectrum is
\[
\boxed{
 -(q-1)^{q(q-1)/2}\oplus
 (q+1)^{q(q-3)/2}\oplus
 (1-\sqrt q)^q\oplus
 (1+\sqrt q)^q .}
\]
After division by $q$, its empirical spectral measure converges to
\[
\boxed{\tfrac12\delta_{-1}+\tfrac12\delta_{+1}.}
\]
Thus the Heisenberg phase removes substantial finite degeneracy but, at maximal all-$q$
regularity, still does not generate a Weyl-law continuum.

The minimum-word theorem of the determinant-isodual code also sharpens.  If $S$ is a
weight-$(q+1)$ word, parity against the opposite determinant coset and the counts
\[
 a=\frac{q(q-1)}2,
 \qquad
 b=\frac{q-1}{2}
\]
of opposite projectivities through one cell and through two compatible cells force equality
in ${m\choose2}\ge m/2$ for every even intersection size $m$.  Consequently $S$ has one
cell in every row and column, and every opposite projectivity meets it in $0$ or $2$ cells.
Thus minimum words are permutation graphs preserving the two $\mathrm{PSL}_2(q)$ colours
of ordered triples.  The standard projective switching automorphism group is
$\mathrm{P}\Sigma L_2(q)$, so field automorphisms occur as genuine minimum histories over
extension fields.  Exhaustive prime-field checks give $12,60,168$ minimum histories at
$q=3,5,7$, while at $q=9$ the Frobenius map $x\mapsto x^3$ is a minimum word although it
is not a projectivity.

A correction is required for the previous sixteen-event magnetic operator.  The spectrum
computed there is exact for the chosen normalized projective representatives, but the phase
$\omega^{B(s,t)}$ is not projectively intrinsic: changing one representative $s\mapsto-s$
changes $\operatorname{tr}A^4$ from $2256$ to $2400$.  This is structural rather than a
numerical defect.  The central element $-I\in\mathrm{Sp}(4,3)$ fixes every projective point
but exchanges its two nonzero vector lifts, so no $\mathrm{Sp}(4,3)$-equivariant section of
projectivization exists.

The intrinsic $q=3$ replacement is the vector lift of the $16$ Segre events to all $32$
vectors $\pm v$, carrying the alternating Heisenberg cocycle
\[
 \psi(v,w)=\tfrac12 B(v,w)=2B(v,w)\pmod3.
\]
Its Hermitian magnetic spectrum is
\[
\boxed{
 -6^6\oplus(-3)^7\oplus(-1)^3\oplus2^6\oplus3^5\oplus6^4\oplus9^1.}
\]
The projective-to-vector fiber has $q-1$ nonzero lifts, so it is literally a two-sheeted
sign cover precisely at $q=3$ among odd $q$.  This supplies a spin/frame-like structural
selector, not yet a derivation of physical spin statistics.

This selector is independent of the fixed-point Bose--Mesner degeneration
\[
 k_2=\frac{q(q+1)(q-3)}4,
 \qquad
 m_1=\frac{(q-3)(q+1)^2}{4},
\]
where both a relation and a primitive idempotent disappear at $q=3$.  The repo's existing
mixing dictionary also has defect
\[
 \sin^2\theta_{23}-\sin^2\theta_W-\sin^2\theta_{12}
 =\frac{q(q-3)}{q^2+q+1},
\]
but that third statement remains explicitly model-dependent.  The older all-$q$ Ollivier
curvature extrapolation is not used as a theorem here; only its exhaustive $q=3$ value is
certified in the older script.

The qutrit/E$_6$ connection gives a concrete finite gauge sector.  The previously derived
$AG(2,3)$ connection satisfies
\[
 F(d_1,d_2)=-\det(d_1,d_2)\pmod3,
\]
so magnetic translations obey the qutrit Heisenberg relation $ZX=\omega XZ$.  The minimal
Harper operator
\[
 H_{\rm H}=X+X^\dagger+Z+Z^\dagger
\]
has exact spectrum
\[
\boxed{-2,\;1-\sqrt3,\;1+\sqrt3.}
\]
With the declared finite Wilson normalization $\sum_p(1-\Re\omega^{F_p})$, the $54$
point/direction-class plaquettes give $81$; this normalization-dependent number is not a
physical coupling prediction.

Finally, couple the intrinsic magnetic operator $H_{\rm mag}$ to the existing causal
Clifford Dirac walk by
\[
 H(p)=\sum_{j=1}^3p_j\alpha_j\otimes I
 +\beta\otimes(m_0I+gH_{\rm mag}).
\]
Clifford anticommutation yields the exact identity
\[
\boxed{H(p)^2=|p|^2I+I\otimes(m_0I+gH_{\rm mag})^2,}
\]
and therefore
\[
\boxed{E_{h,\pm}(p)=\pm\sqrt{|p|^2+(m_0+gh)^2}.}
\]
This is an exact relativistic-form dispersion in lattice units.  The parameters $m_0,g$,
physical $c$, particle assignments, and curved-spacetime dynamics remain open.

A useful matter-sector falsifier also emerges.  Every allowed E$_6$ cubic lift has
$(t_0,t_1,t_2)=(k,k+\lambda,k+2\lambda)$ and hence total $\mathbb Z_3$ charge zero, but
each forbidden firewall fiber $(0,1,2)$ is neutral as well.  Thus simple $\mathbb Z_3$
charge conservation cannot explain the $36/9$ cubic selection.  The selector must involve
richer incidence, holonomy, $L_\infty$, CE2/metaplectic, or chirality data.

\paragraph{Cover/F$_4$ boundary.}
The selected $13$-cover still has exact stabilizer data
$1152\to576$ with kernel $2$ and orbit split $1+12$.  The direct GAP program now requests
an explicit moving-$12$ conjugator and the committed consumer will compose it with the
frozen Latin--F$_4$ short-root-pair map.  No cover-to-root object dictionary is asserted
until that queued action-level witness is emitted.

\paragraph{Evidence boundary.}
The finite algebra, code, spectrum, gauge-curvature, and Dirac identities above are kept
separate from physical semantics.  They do not yet derive the Standard Model spectrum,
physical masses, the speed of light, a gauge coupling, gravity, or a continuum spacetime.
